% secV6_series.tex -- Section 5. CONTAINER (27 Oct 2026): text to be written by the author.
% Word budget: 700 words (JPART count: running text, footnotes and citations as printed).
% Outline, numbers and sources: referee_EE_2026-10/ESQUELETO.md, section 5. The archived text
% (_arquivo_texto_IA_2026-10-27/) is not a source.

\section{The published number}\label{sec:series}\label{sec:stock}

The published waiting list is an administrative count whose definition and observation date have changed over time. In 2011, the inclusion of Minigrupo II altered the population covered by childcare statistics \citep{sme2011nota_reorg}. A July ordinance allowed families already referred elsewhere to retain a centre-specific request \citep{sme2011portaria3478}; an October publication identified these cases as \emph{demanda preferencial} \citep{sme2011nota_pref}. In March 2012, SME distinguished the stock of requests carried over from the previous year from new registrations and identified cases pending documentation \citep{sme2012nota}. From September 2019, releases separated requests with and without a specific-centre preference; from September 2020, they also distinguished waiting times up to and above 30 days. Publication changed from quarterly to bimonthly in 2023. Accordingly, levels cannot be treated as a consistently defined time series across 2011 or 2019 (Table~\ref{tab:series}). The distinction remains consequential: a December 2020 municipal announcement described the narrower waiting list as cleared while acknowledging 540 children waiting for a specific centre \citep{saopaulo2020fila}. The published list itself was not zero. The 2025--2028 municipal goals programme uses a still narrower measure: requests without a specific-centre preference pending for more than 30 days \citep{sme_metas2025}.

Even within periods of comparable definitions, the date of observation substantially changes the count. December has the year's lowest published waiting-list count in 17 of the 19 years from 2007 to 2025. Between 2013 and 2017, the December count fell from 96,696 to 44,092, a 54.4\% reduction; the corresponding changes in March, June, and September were only 14.9\%, 18.1\%, and 15.7\% (Table~\ref{tab:series}). The within-year decline from September to December rose from 38.4\% in 2013 to 87.2\% in 2019. The annual mean of the four quarterly counts declined by 23.9\% between 2013 and 2017, from 121,087 to 92,158 (Table~\ref{tab:series}). Thus, both the measured level and its apparent trend depend on the month chosen.

Individual administrative records allow us to investigate the mechanism for 2013--2018. We reconstruct the stock of unresolved first requests at each quarter-end using dated entries and exits from the request record. Across those quarters, reconstructed counts are 87--96\% of the corresponding published counts (Table~\ref{tab:recon}). They exhibit the same seasonal pattern: from 2013 to 2017, the reconstructed December stock fell by 53\%, against declines of 16--19\% in the other quarters; Table~\ref{tab:decomp} reports the December and September levels. The December decline is therefore not confined to the publication series, although the differences in levels remain unexplained.

A stock--flow decomposition identifies referrals as the principal source of the year-end reduction. Over 2013--2018, 97.4\% of recorded exits between September and December are referrals, compared with 2.3\% attributable to ageing out and 0.3\% to recorded withdrawals (Table~\ref{tab:decomp}). For 2013--2017, referrals during this interval average 94,262 per year, equivalent to 71.2\% of the reconstructed September stock; approximately 32,150 new requests also enter during the interval. These movements are consistent with the main allocation round described in Section~\ref{sec:institutions}, rather than a year-end purge of registrations. The decomposition is unavailable after 2018 and does not establish the mechanism in later years.

The publication's extraction date further affects what a ``December'' observation captures. For 2014--2018, at least 99.2\% of referrals occurring in December had already taken place by the corresponding extraction date (Table~\ref{tab:decomp}). These snapshots therefore follow the December wave of referrals, though they precede the subsequent January--February wave. The extraction convention is not uniform: of 18 December editions with a printed date, ten use December 31, while other dates range from November 30 to December 28 (Table~\ref{tab:dates}). These dates establish variation, not its administrative cause.

The choice carries directly into assessments of Goal~5. Applying the 75\% threshold used in the municipality's 2023 review to the share $E^B_t/(E^B_t+L_t)$ yields a first crossing in 2015 with December counts, in 2018 with September counts, and in 2017 using the annual mean of published counts (Table~\ref{tab:series}). This threshold is not the law's requirement to serve 100\% of registered demand. The December count, taken after the year's principal referral wave, consequently makes the review's criterion appear to be reached earlier (Figure~\ref{fig:goal}). It is an observable mechanism behind the reference-date sensitivity documented in Section~\ref{sec:goal}, not evidence that the statutory goal was met.
